\documentclass[reqno]{amsart}
\usepackage[T1]{fontenc}
\usepackage{amsmath,amssymb,amsthm}
\usepackage{xr-hyper}
\usepackage[hidelinks]{hyperref}
\externaldocument[sd-]{20261125T000000--slope-method__workbook}
\externaldocument[at-]{20261126T030100--a-twist-for-each-prime__workbook}
\externaldocument[ds-]{20261125T004000--denominators-and-sha__workbook}
\externaldocument[es-]{20261125T014500--eisenstein-congruences-and-selmer-classes__workbook}
\externaldocument[da-]{20261126T030300--the-data__workbook}
\externaldocument[col-]{20261126T030800--the-cost-of-one-lift__workbook}
\externaldocument[dp-]{20261126T030900--denominators-of-projections-to-lifts__workbook}
\externaldocument[n3-]{20261126T051800--the-local-loss-and-the-sparse-places__workbook}
\externaldocument[tc-]{20261126T040200--theorem-c__workbook}
\externaldocument[w1-]{20261126T040300--roots-of-unity-for-many-polynomials__workbook}
\externaldocument[ob-]{20261126T040100--spaces-forms-and-twists__workbook}
\externaldocument[l5-]{20261126T051100--test-forms-and-the-height-of-a-hecke-ratio__workbook}

% Chapter number; \ref{chap} gives it.
\newcounter{chap}
% Chapter numbers, generated from 20261125T012000--grh-workbook__guide.org
% by scripts/make-chapters.py.  Do not edit: change the order of the
% index and rerun the script (./preview-site.sh and CI do this).
%
% Usage in a sheet, after \newcounter{chap}:
%   % Chapter numbers, generated from 20261125T012000--grh-workbook__guide.org
% by scripts/make-chapters.py.  Do not edit: change the order of the
% index and rerun the script (./preview-site.sh and CI do this).
%
% Usage in a sheet, after \newcounter{chap}:
%   % Chapter numbers, generated from 20261125T012000--grh-workbook__guide.org
% by scripts/make-chapters.py.  Do not edit: change the order of the
% index and rerun the script (./preview-site.sh and CI do this).
%
% Usage in a sheet, after \newcounter{chap}:
%   \input{chapters}\chapnumber{<sheet ID>}
\makeatletter
\providecommand\chapters@set[2]{\@namedef{chapters@#1}{#2}}
\chapters@set{20261125T001000}{1}% zeros-to-ramp-sums
\chapters@set{20261125T001100}{2}% inert-primes-and-estimate-b
\chapters@set{20261125T005000}{3}% galois-fourier-and-averaging
\chapters@set{20261125T005500}{4}% columns-from-a-sphere
\chapters@set{20261125T000000}{5}% slope-method
\chapters@set{20261126T030100}{6}% a-twist-for-each-prime
\chapters@set{20261125T003000}{7}% the-chain-and-the-square-root
\chapters@set{20261125T010000}{8}% rank-lemma-and-zero-estimate
\chapters@set{20261125T010500}{9}% poincare-series-and-near-orthogonality
\chapters@set{20261124T223000}{10}% elliptic-curves-for-the-proof
\chapters@set{20261125T009500}{11}% fock-model-and-theta-intertwiners
\chapters@set{20261124T223100}{12}% complex-multiplication-and-inert-reduction
\chapters@set{20261125T008000}{13}% one-cm-determinant
\chapters@set{20261124T214500}{14}% supersingular-theta-divisibility
\chapters@set{20261124T214600}{15}% the-isotropic-plane
\chapters@set{20261125T013700}{16}% pairings-on-a-formal-group-of-height-two
\chapters@set{20261125T013500}{17}% hecke-nondegeneracy-on-the-tree
\chapters@set{20261126T030200}{18}% the-bound-is-attained
\chapters@set{20261125T009000}{19}% kronecker-limit-and-chowla-selberg
\chapters@set{20261125T002000}{20}% definite-theta-lifts
\chapters@set{20261125T015500}{21}% siegel-weil-in-the-smallest-case
\chapters@set{20261125T002100}{22}% definite-theta-lifts-and-norms
\chapters@set{20261125T016500}{23}% hecke-operators-and-theta-lifts
\chapters@set{20261125T016700}{24}% divisibility-along-the-chain
\chapters@set{20261125T004000}{25}% denominators-and-sha
\chapters@set{20261125T004100}{26}% tate-twists-and-sha
\chapters@set{20261126T020430}{27}% the-minus-part-by-gauss-sums
\chapters@set{20261125T014100}{28}% congruence-modules
\chapters@set{20261125T014500}{29}% eisenstein-congruences-and-selmer-classes
\chapters@set{20261125T014700}{30}% pseudocharacters-and-the-length-bound
\chapters@set{20261125T011000}{31}% galois-representations-black-box
\chapters@set{20261126T030300}{32}% the-data
\chapters@set{20261126T020210}{33}% bessel-poincare-series-and-their-gram-operators
\chapters@set{20261126T030400}{34}% the-seesaw-identity
\chapters@set{20261126T030500}{35}% divisibility-of-the-composite-lift
\chapters@set{20261126T030600}{36}% theorem-a-short-forms
\chapters@set{20261126T030700}{37}% the-determinant-of-a-chain
\chapters@set{20261126T030800}{38}% the-cost-of-one-lift
\chapters@set{20261126T030900}{39}% denominators-of-projections-to-lifts
\chapters@set{20261126T031000}{40}% twists-with-few-congruences
\chapters@set{20261126T031100}{41}% theorem-c-and-the-main-inequality
\providecommand\chapnumber[1]{%
  \@ifundefined{chapters@#1}%
    {\PackageError{chapters}{Sheet #1 is not listed in the index}%
      {List it in 20261125T012000--grh-workbook__guide.org\MessageBreak
       and run scripts/make-chapters.py.}}%
    {\edef\thechap{\@nameuse{chapters@#1}}}}
\makeatother
\chapnumber{<sheet ID>}
\makeatletter
\providecommand\chapters@set[2]{\@namedef{chapters@#1}{#2}}
\chapters@set{20261125T001000}{1}% zeros-to-ramp-sums
\chapters@set{20261125T001100}{2}% inert-primes-and-estimate-b
\chapters@set{20261125T005000}{3}% galois-fourier-and-averaging
\chapters@set{20261125T005500}{4}% columns-from-a-sphere
\chapters@set{20261125T000000}{5}% slope-method
\chapters@set{20261126T030100}{6}% a-twist-for-each-prime
\chapters@set{20261125T003000}{7}% the-chain-and-the-square-root
\chapters@set{20261125T010000}{8}% rank-lemma-and-zero-estimate
\chapters@set{20261125T010500}{9}% poincare-series-and-near-orthogonality
\chapters@set{20261124T223000}{10}% elliptic-curves-for-the-proof
\chapters@set{20261125T009500}{11}% fock-model-and-theta-intertwiners
\chapters@set{20261124T223100}{12}% complex-multiplication-and-inert-reduction
\chapters@set{20261125T008000}{13}% one-cm-determinant
\chapters@set{20261124T214500}{14}% supersingular-theta-divisibility
\chapters@set{20261124T214600}{15}% the-isotropic-plane
\chapters@set{20261125T013700}{16}% pairings-on-a-formal-group-of-height-two
\chapters@set{20261125T013500}{17}% hecke-nondegeneracy-on-the-tree
\chapters@set{20261126T030200}{18}% the-bound-is-attained
\chapters@set{20261125T009000}{19}% kronecker-limit-and-chowla-selberg
\chapters@set{20261125T002000}{20}% definite-theta-lifts
\chapters@set{20261125T015500}{21}% siegel-weil-in-the-smallest-case
\chapters@set{20261125T002100}{22}% definite-theta-lifts-and-norms
\chapters@set{20261125T016500}{23}% hecke-operators-and-theta-lifts
\chapters@set{20261125T016700}{24}% divisibility-along-the-chain
\chapters@set{20261125T004000}{25}% denominators-and-sha
\chapters@set{20261125T004100}{26}% tate-twists-and-sha
\chapters@set{20261126T020430}{27}% the-minus-part-by-gauss-sums
\chapters@set{20261125T014100}{28}% congruence-modules
\chapters@set{20261125T014500}{29}% eisenstein-congruences-and-selmer-classes
\chapters@set{20261125T014700}{30}% pseudocharacters-and-the-length-bound
\chapters@set{20261125T011000}{31}% galois-representations-black-box
\chapters@set{20261126T030300}{32}% the-data
\chapters@set{20261126T020210}{33}% bessel-poincare-series-and-their-gram-operators
\chapters@set{20261126T030400}{34}% the-seesaw-identity
\chapters@set{20261126T030500}{35}% divisibility-of-the-composite-lift
\chapters@set{20261126T030600}{36}% theorem-a-short-forms
\chapters@set{20261126T030700}{37}% the-determinant-of-a-chain
\chapters@set{20261126T030800}{38}% the-cost-of-one-lift
\chapters@set{20261126T030900}{39}% denominators-of-projections-to-lifts
\chapters@set{20261126T031000}{40}% twists-with-few-congruences
\chapters@set{20261126T031100}{41}% theorem-c-and-the-main-inequality
\providecommand\chapnumber[1]{%
  \@ifundefined{chapters@#1}%
    {\PackageError{chapters}{Sheet #1 is not listed in the index}%
      {List it in 20261125T012000--grh-workbook__guide.org\MessageBreak
       and run scripts/make-chapters.py.}}%
    {\edef\thechap{\@nameuse{chapters@#1}}}}
\makeatother
\chapnumber{<sheet ID>}
\makeatletter
\providecommand\chapters@set[2]{\@namedef{chapters@#1}{#2}}
\chapters@set{20261125T001000}{1}% zeros-to-ramp-sums
\chapters@set{20261125T001100}{2}% inert-primes-and-estimate-b
\chapters@set{20261125T005000}{3}% galois-fourier-and-averaging
\chapters@set{20261125T005500}{4}% columns-from-a-sphere
\chapters@set{20261125T000000}{5}% slope-method
\chapters@set{20261126T030100}{6}% a-twist-for-each-prime
\chapters@set{20261125T003000}{7}% the-chain-and-the-square-root
\chapters@set{20261125T010000}{8}% rank-lemma-and-zero-estimate
\chapters@set{20261125T010500}{9}% poincare-series-and-near-orthogonality
\chapters@set{20261124T223000}{10}% elliptic-curves-for-the-proof
\chapters@set{20261125T009500}{11}% fock-model-and-theta-intertwiners
\chapters@set{20261124T223100}{12}% complex-multiplication-and-inert-reduction
\chapters@set{20261125T008000}{13}% one-cm-determinant
\chapters@set{20261124T214500}{14}% supersingular-theta-divisibility
\chapters@set{20261124T214600}{15}% the-isotropic-plane
\chapters@set{20261125T013700}{16}% pairings-on-a-formal-group-of-height-two
\chapters@set{20261125T013500}{17}% hecke-nondegeneracy-on-the-tree
\chapters@set{20261126T030200}{18}% the-bound-is-attained
\chapters@set{20261125T009000}{19}% kronecker-limit-and-chowla-selberg
\chapters@set{20261125T002000}{20}% definite-theta-lifts
\chapters@set{20261125T015500}{21}% siegel-weil-in-the-smallest-case
\chapters@set{20261125T002100}{22}% definite-theta-lifts-and-norms
\chapters@set{20261125T016500}{23}% hecke-operators-and-theta-lifts
\chapters@set{20261125T016700}{24}% divisibility-along-the-chain
\chapters@set{20261125T004000}{25}% denominators-and-sha
\chapters@set{20261125T004100}{26}% tate-twists-and-sha
\chapters@set{20261126T020430}{27}% the-minus-part-by-gauss-sums
\chapters@set{20261125T014100}{28}% congruence-modules
\chapters@set{20261125T014500}{29}% eisenstein-congruences-and-selmer-classes
\chapters@set{20261125T014700}{30}% pseudocharacters-and-the-length-bound
\chapters@set{20261125T011000}{31}% galois-representations-black-box
\chapters@set{20261126T030300}{32}% the-data
\chapters@set{20261126T020210}{33}% bessel-poincare-series-and-their-gram-operators
\chapters@set{20261126T030400}{34}% the-seesaw-identity
\chapters@set{20261126T030500}{35}% divisibility-of-the-composite-lift
\chapters@set{20261126T030600}{36}% theorem-a-short-forms
\chapters@set{20261126T030700}{37}% the-determinant-of-a-chain
\chapters@set{20261126T030800}{38}% the-cost-of-one-lift
\chapters@set{20261126T030900}{39}% denominators-of-projections-to-lifts
\chapters@set{20261126T031000}{40}% twists-with-few-congruences
\chapters@set{20261126T031100}{41}% theorem-c-and-the-main-inequality
\providecommand\chapnumber[1]{%
  \@ifundefined{chapters@#1}%
    {\PackageError{chapters}{Sheet #1 is not listed in the index}%
      {List it in 20261125T012000--grh-workbook__guide.org\MessageBreak
       and run scripts/make-chapters.py.}}%
    {\edef\thechap{\@nameuse{chapters@#1}}}}
\makeatother
\chapnumber{20261126T031000}
\renewcommand{\thesection}{\thechap.\arabic{section}}

\theoremstyle{definition}
\newtheorem{problem}{Problem}[section]
\theoremstyle{plain}
\newtheorem{theorem}{Theorem}[section]

\newcommand{\Z}{\mathbb{Z}}
\newcommand{\Q}{\mathbb{Q}}
\newcommand{\C}{\mathbb{C}}
\newcommand{\Ql}{\bar{\mathbb{Q}}_\ell}
\newcommand{\Zl}{\bar{\mathbb{Z}}_\ell}
\newcommand{\Sel}{\mathrm{Sel}}
\newcommand{\val}{\operatorname{val}}
\newcommand{\Oc}{\mathcal{O}}
\newcommand{\St}{\mathrm{St}}
\newcommand{\bd}{\mathrm{bd}}
\newcommand{\TS}{T_{\mathrm{Sel}}}
\newcommand{\ctor}{c_{\mathrm{tor}}}
\newcommand{\cl}{\mathrm{cl}}
\newcommand{\loc}{\mathrm{loc}}

\setlength{\emergencystretch}{3em}

\title{Twists with few congruences}
\author{}
\date{November 26, 2026}

\begin{document}
\maketitle
\refstepcounter{chap}\label{chap}

\paragraph{Prerequisites.}
Valuations on \(\Ql\) \cite[Ch.~II]{Ne}; roots of unity of prime power order \cite[Ch.~1--2]{Wa}; inertia groups and tame inertia at a finite place.  Chapter~\ref{w1-chap}, \emph{Roots of unity for many polynomials}: Theorem~\ref{w1-thm:roots} and Problems~\ref{w1-pr:powers} and~\ref{w1-pr:powers-far}.  Chapter~\ref{dp-chap}, \emph{Denominators of projections to lifts}: Section~\ref{dp-sec:statement} with its definitions (coefficient rings and normalized lengths, auxiliary twists, main-summand-like representations, Selmer bounds, the numbers \(\mathrm{tw}_\ell\) and \(\cl_\ell\)) and Theorem~\ref{dp-thm:denominators}.  The data of Chapter~\ref{da-chap}; twists and critical pairs (Chapter~\ref{ob-chap}, Section~\ref{ob-sec:twists} and Definition~\ref{ob-def:pair}); heights of families (Section~\ref{sd-sec:heights}) and the twist at each prime of Chapter~\ref{at-chap}; regular families (Definitions~\ref{l5-def:small-ratios} and~\ref{l5-def:regular}).  As background: Stickelberger elements (Chapter~\ref{ds-chap}) and Selmer groups (Chapter~\ref{es-chap}).

\medskip
Chapter~\ref{dp-chap} bounds the denominator of a projection to lifts through the Selmer groups of twisted Galois representations, and the cost of comparing a lift at two twists involves ratios of Hecke \(L\)-values (Chapter~\ref{col-chap}).  For each prime \(\ell\) and each \(\iota \mid \ell\), can one twist bound the Selmer lengths for a prescribed finite family satisfying the hypotheses of Theorem~\ref{thm:twisting} and the valuations of the ratios for every ramified critical pair and both signs?  The proof reduces both to values of polynomials at points \((\zeta_w-1)_w\), with \(\zeta_w\) roots of unity of \(\ell\)-power order.  The lemma that chooses the roots of unity is proved in Chapter~\ref{w1-chap}; Section~\ref{sec:twisting} states the polynomial input, deduces the existence of small twists (Theorem~\ref{thm:small}) and states the twisting theorem; Section~\ref{sec:regular} provides the regular families used in Chapters~\ref{col-chap} and~\ref{tc-chap}.

Throughout, \(\ell\) is a prime and \(\val\) is the valuation of \(\Ql\) with \(\val(\ell) = 1\) and \(\val(0) = \infty\), as in Chapter~\ref{w1-chap}; \(\val_+ := \max(\val, 0)\), and \(\val_\ell\) is the \(\ell\)-adic valuation of integers.  Coefficient rings \(\Oc\) and normalized lengths are those of Section~\ref{dp-sec:statement}.

\section{The twisting theorem}\label{sec:twisting}

Fix a CM field \(E\subseteq\bar\Q\), with \(G_E := \mathrm{Gal}(\bar\Q/E)\) and \(g^c := cgc\) for the complex conjugation \(c\) of Section~\ref{dp-sec:statement}, a finite set \(\Sigma_0\) of finite places of \(E\) containing those above \(\ell\), and an element \(g_0\) of inertia at a place of \(\Sigma_0\) prime to \(\ell\).  The proof works with finitely many objects of two kinds, Selmer data and Stickelberger systems; we describe them only as far as this section uses them.  Chapter~\ref{es-chap} gives examples of Selmer classes arising from congruences, and Chapter~\ref{ds-chap} treats the classical cyclotomic Stickelberger element.

The Selmer data \(\TS\) have a common coefficient ring \(\Oc\).  Each is a lattice of rank \(\operatorname{rk} \TS\) with continuous action of \(G_E\) unramified outside a set \(\Sigma(\TS) \subseteq \Sigma_0\) containing the places above \(\ell\), with local cocycle conditions above \(\ell\) and a distinguished inertial element \(\gamma(\TS)\).  Write \(\bd(\TS)\) for the normalized sum of its global and local terms, which include the valuations of \(\det(\TS(g) - 1)\) for \(g = g_0\) and \(g = (g_0^c)^{-1}\).  For a finite set \(Q\) of places and a character \(\chi\) of \(G_E\), a bound \(\mathrm{len}\,\Sel_Q(\TS\chi) \le d\) means that for every sufficiently large coefficient ring \(\Oc' \supseteq \Oc\) there is an \(\Oc'\)-valued model of \(\chi\) for which the following module has normalized length at most \(d\): the locally constant cocycles with coefficients in the extended twisted lattice tensored with \(\mathrm{Frac}(\Oc')/\Oc'\) that are unramified outside \(\Sigma(\TS) \cup Q \cup \bar Q\), satisfy the local relations and vanish at \(\gamma(\TS)\).  A Stickelberger system \(\St\) has elements \(\St_Q \in \Zl[\mathrm{Cl}_Q]\), where \(\mathrm{Cl}_Q\) is the norm-one idele class group of level \(Q\), and a finite number \(\bd(\St)\) attached to the system, whose definition we shall not need.  A finite set \(Q\) of places is \emph{auxiliary} for \(\Sigma_0\) if its places have prime norm, and \(Q\), its conjugate \(\bar Q\) and \(\Sigma_0\) are pairwise disjoint.

The polynomial input below is read from the proof; it follows by descent, which we do not treat.  It holds for any finite \(\Sigma_0\) with \(g_0\) inertial at a place of \(\Sigma_0\) prime to \(\ell\), subject to the hypotheses on the data, systems and \(\ctor\) below.  We use it with \(\Sigma_0\) containing the places above \(\ell\).  Fix a natural number \(\ctor\), a \emph{torsion level}: for every datum, \(\bd(\TS) + 1 < \ctor\) and \(2\operatorname{rk} \TS\,\mathrm{ord}_{\Oc}(\ell^{\ctor}) + 1 < \#k_{\Oc}\), where \(k_{\Oc}\) is the residue field and \(\mathrm{ord}_{\Oc}\) gives a uniformizer the value \(1\).  There is an integer \(s_0 \ge 1\) such that for every rational \(\rho > 0\) with \(s_0\rho < 1\) and every sufficiently large \(h_1\) there are:
\begin{enumerate}
\item an auxiliary set \(Q\) of at most \(s_0/2\) places, and for each \(w \in Q\) a locally constant homomorphism \(\lg_w\colon G_E \to \Z/\ell^{h_1}\), onto on inertia at \(w\), coming by \(y \mapsto y/\bar y\) from a finite idele class character of tame level \(\{w\}\) trivial above \(\ell\); thus it is anticyclotomic, unramified outside \(w\) and \(\bar w\), and trivial on the decomposition groups above \(\ell\);
\item for each datum and each system a polynomial \(F\) over \(\Zl\) in variables \(X_w\), \(w \in Q\), and a point with coordinates of valuation at least \(\rho\) at which \(\val(F) \le \bd + s_0\rho\), where \(\bd\) is \(\bd(\TS)\) or \(\bd(\St)\);
\item for \(\ell^{h_1}\)-th roots of unity \(\zeta_w\), with \(\chi_\zeta := \prod_{w \in Q} \zeta_w^{\lg_w}\) and \(x := (\zeta_w - 1)_w\): for a system, \(F(x) = \chi_\zeta(\St_Q)\); for a datum with \(\val(F(x)) < \bd(\TS) + 1\), \(\mathrm{len}\,\Sel_Q(\TS\chi_\zeta) \le \val(F(x))\).
\end{enumerate}

\begin{problem}\label{pr:psi-power}
For \(\chi_\zeta = \prod_{w} \zeta_w^{\lg_w}\) and \(n \ge 1\), \(\chi_\zeta^n = \prod_{w} (\zeta_w^n)^{\lg_w}\), and \(\chi_\zeta^{\ell^{h_1}} = 1\).  For \(w \in Q\) some \(g\) in an inertia group at \(w\) has \(\chi_\zeta(g) = \zeta_w\).
\end{problem}

\begin{problem}\label{pr:inertia-sum}
Let \(0 < \rho < 1\) and \(m \ge 1\), and let the \(\zeta_w\), \(w \in Q\), be \(\ell^{h_1}\)-th roots of unity with \(\val(\zeta_w - 1) \le \rho/m\).  Then \(\val\bigl(\chi_\zeta^n(s_{w'}) - 1\bigr) \le \rho\) for every \(1 \le n \le m\), every \(w' \in Q \cup \bar Q\) and every generator \(s_{w'}\) of tame inertia at a place of \(\bar\Q\) above \(w'\).  Hence \(\sum_{w' \in Q \cup \bar Q}\val\bigl(\chi_\zeta^n(s_{w'}) - 1\bigr) \le 2\,\#Q\,\rho\).
\end{problem}

\begin{problem}\label{pr:twist-from-roots}
Let \(0 < \varepsilon < 1\) and \(m \ge 1\).  There are auxiliary data as in the input and a character \(\chi_\zeta = \prod_{w} \zeta_w^{\lg_w}\) such that, for all \(1 \le n \le m\), \(\mathrm{len}\,\Sel_Q(\TS\chi_\zeta^n) \le \bd(\TS) + \varepsilon\) for every datum, \(\val(\chi_\zeta^n(\St_Q)) \le \bd(\St) + \varepsilon\) for every system, \(\chi_\zeta^n\) is ramified at every place of \(Q\), and \(\sum_{w' \in Q \cup \bar Q}\val\bigl(\chi_\zeta^n(s_{w'}) - 1\bigr) \le \varepsilon\) for any generators \(s_{w'}\) of tame inertia.
\end{problem}

\begin{theorem}[small twists]\label{thm:small}
Let \(\Sigma_0\), \(g_0\), finitely many Selmer data over one coefficient ring \(\Oc\), finitely many Stickelberger systems and a torsion level \(\ctor\) be as above, and let \(m \ge 1\) be an integer and \(\varepsilon > 0\) real.  There are a finite set \(Q\) of places, auxiliary for \(\Sigma_0\), and a character \(\chi\colon G_E \to \Ql^\times\) such that for every \(1 \le n \le m\):
\begin{enumerate}
\item \(\chi^n\) is locally constant, of \(\ell\)-power order, comes by \(y \mapsto y/\bar y\) from a finite idele class character of tame level \(Q\) trivial above \(\ell\), and is ramified at every place of \(Q\);
\item \(\mathrm{len}\,\Sel_Q(\TS\chi^n) \le \bd(\TS) + \varepsilon\) for every datum, and \(\val(\chi^n(\St_Q)) \le \bd(\St) + \varepsilon\) for every system;
\item \(\sum_{w' \in Q \cup \bar Q}\val\bigl(\chi^n(s_{w'}) - 1\bigr) \le \varepsilon\) for any generators \(s_{w'}\) of tame inertia at any places of \(\bar\Q\) above the \(w'\).
\end{enumerate}
\end{theorem}

For \(\varepsilon < 1\), Problems~\ref{pr:psi-power}--\ref{pr:twist-from-roots} deduce Theorem~\ref{thm:small} from the polynomial input, with \(\chi = \chi_\zeta\); for \(\varepsilon \ge 1\) it follows from the case \(\varepsilon = 1/2\).  The proof proceeds in the same way.

For the data of Chapter~\ref{da-chap} the proof applies Theorem~\ref{thm:small}, with \(\varepsilon = 2^{-\ell}\), to Selmer data on lattices in given representations and to one Stickelberger system for each ramified critical pair and sign.  We use the notation of Chapters~\ref{da-chap} and~\ref{ob-chap} (\(f\), \(r\), \(t\), \(m\), \(b\), \(p_0\), \(p_1\), \(M = p_0p_1\), \(N\), the chosen places \(\mathcal M\), the integer \(\mathcal N\), \(m\)-fold twists \(\vartheta\) with their sets \(Q_\vartheta\)), the critical pairs \(P = (\nu, z)\), their ratios \(R_P(\vartheta, \epsilon)\), \(\epsilon = \pm1\), and ramified pairs of Definition~\ref{ob-def:pair}, and the number of Definition~\ref{dp-def:clearing-bound},
\[
  \mathrm{tw}_\ell := (4f + 3)\, r\, \val_\ell(\mathcal N) + 2^{-\ell}.
\]
For a twist \(\vartheta\) and \(\iota\colon \C \simeq \Ql\), \(\hat\vartheta\) is the inverse avatar of Definition~\ref{dp-def:aux-twist}.

\begin{theorem}[twisting theorem]\label{thm:twisting}
Let \(\iota\colon \C \simeq \Ql\), and let \(R_1, \dots, R_n\) (\(n \ge 0\)) be \(\ell\)-adic representations of \(G_E\) such that each \(R_i\) is main-summand-like for \((t'_i, \mathrm{HT}_i)\), for some integer \(t'_i \ge 0\) and some family \(\mathrm{HT}_i\) (Definition~\ref{dp-def:main-summand-like}).  There is an \(m\)-fold twist \(\vartheta\) such that for every \(1 \le j \le m\):
\begin{enumerate}
\item the twist \(\vartheta^j\) is auxiliary with ramification at most \(\mathrm{tw}_\ell\) (Definition~\ref{dp-def:aux-twist});
\item for each \(i\), the twist \(\widehat{\vartheta^j}R_i\) has Selmer bound \(\mathrm{tw}_\ell\) (Definition~\ref{dp-def:selmer-bound});
\item for each ramified critical pair \(P\) and each sign \(\epsilon\), \(\val_+\bigl(\iota(R_P(\vartheta^j, \epsilon))\bigr) \le \mathrm{tw}_\ell\).
\end{enumerate}
\end{theorem}

\emph{Status}: read from the proof \cite{GL}; not proved here.  The proof applies Theorem~\ref{thm:small} with \(\Sigma_0\) the places dividing \(\ell MN\).  The hypothesis on the \(R_i\) imposes, at each place of \(\bar\Q\) above a chosen place away from \(\ell\), the determinant bounds for the distinguished tame generator; it does not include the trace identity of Definition~\ref{dp-def:main-summand}.  The orders of finite general linear groups enter \(\mathrm{tw}_\ell\) through \(\mathcal N\); the next problem is the smallest instance.

\begin{problem}\label{pr:group-order}
\(\#\mathrm{GL}_3(\mathbb F_7) = 33784128 = 2^6\cdot 3^4\cdot 7^3\cdot 19\), and \(\val_3(7^e - 1) = 1, 1, 2\) for \(e = 1, 2, 3\), with sum \(4 = \val_3(\#\mathrm{GL}_3(\mathbb F_7))\).
\end{problem}

\begin{problem}\label{pr:vl-log}
For an integer \(n \ge 1\) and a finite set \(P_0\) of primes, \(\sum_{\ell \in P_0} \val_\ell(n)\log \ell \le \log n\).  Moreover \(\sum_\ell 2^{-\ell}\log \ell = 0.3776\ldots\), the sum over all primes.  Hence \(\sum_{\ell \in P_0} \mathrm{tw}_\ell \log \ell \le (4f + 3)\, r \log \mathcal N + 1\).
\end{problem}

\section{Regular families}\label{sec:regular}

As in Chapter~\ref{at-chap}, the twist may depend on the prime: it is chosen for every \(\iota \mid \ell\), and families \(y = (y_\iota)\) are measured by their heights (Section~\ref{sd-sec:heights}).  Regular families and small ratios are those of Definitions~\ref{l5-def:regular} and~\ref{l5-def:small-ratios}; the denominators are measured against the budget \(B = U/(mk)\), \(U = fa\log^3 b\).

\begin{problem}\label{pr:height}
Let \(e_\ell \ge 0\) be reals with \(\sum_{\ell \in P_0} e_\ell\log \ell \le c\) for every finite set \(P_0\) of primes.  If \(\lVert \iota(y_\iota)\rVert \le \ell^{e_\ell}\) for every \(\iota \mid \ell\), then \(y\) has height at most \(c\).
\end{problem}

\begin{problem}\label{pr:inverse}
If \(y \in \C\) and \(\val_+(\iota(y)) \le c\) for a real \(c \ge 0\), then \(y \ne 0\) and \(\lVert\iota(y^{-1})\rVert \le \ell^{c}\).
\end{problem}

\begin{problem}\label{pr:coprime}
If \(Q\) is auxiliary for a set \(\Sigma_0\) containing the places above \(\ell\), then \(\ell\) does not divide the product of the norms of the places of \(Q\).
\end{problem}

The clearing bound \(\cl_\ell\) of Definition~\ref{dp-def:clearing-bound} contains the local loss \(\loc_\ell = 4f\,\mathrm{defect}_\ell\) of Definition~\ref{n3-def:loss}, in which \(\mathrm{Int}_b\) is the product of the positive integers in \((b - 30(\Delta/2)^3, b + 30(\Delta/2)^3]\) and \(\Delta/2 = r^8\).  For the data, \(m \ge 1\), \(k \ge 1\) and \(b > 1\).

\begin{problem}\label{pr:loss-sum}
For every finite set \(P_0\) of primes, \(\sum_{\ell \in P_0} \mathrm{defect}_\ell \log \ell \le 10^6 \frac{m}{k}\bigl(b\log M + \log\bigl((2b)!\,(b + 30(\Delta/2)^3)^{60(\Delta/2)^6}\bigr)\bigr)\).
\end{problem}

By Theorem~\ref{n3-thm:local-loss}, \(\sum_{\ell \in P_0} \loc_\ell\log \ell \le 4\cdot 10^6 e^{180} f\frac{m}{k}\, b\log b\) for every finite set \(P_0\) of primes.

\begin{problem}\label{pr:budget}
\(4\cdot 10^6 e^{180} f\frac{m}{k}\, b \log b \le B\) if and only if \(4\cdot 10^6 e^{180}\, m^2 \le x\log^2 b\).  In particular this holds if \(m^2 \le 4x\) and \(\log^2 b \ge 16\cdot 10^6 e^{180}\).
\end{problem}

Problem~\ref{pr:budget} compares the sum of the local losses with the budget; the data have \(x \le m^2 \le 4x\).  The other terms of \(\cl_\ell\) and the class number of \(U(1)\) are summed in Problem~\ref{dp-pr:budget-sum}, and the total is at most \(B\).

\begin{theorem}\label{thm:regular}
For all data \(D\) there is a regular family \(\theta = (\theta_\iota)\) of \(m\)-fold twists.
\end{theorem}

\emph{Status}: read from the proof \cite{GL}.  The next two problems give the clause on small ratios and one step of the clause on denominators.

\begin{problem}\label{pr:assembly}
Assume \(\sum_{\ell \in P_0} \mathrm{tw}_\ell\log \ell \le x_1\) for every finite set \(P_0\) of primes.  If each \(\theta_\iota\) satisfies the conclusions of Theorem~\ref{thm:twisting} at \(\iota\), then \(\theta\) has small ratios.
\end{problem}

\begin{problem}\label{pr:class-number}
Let \(C_1 \ge 1\) be an integer, \(y \in \C\) and \(c\) real.  If \(\lVert\iota(C_1y)\rVert \le e^{c}\), then \(\lVert\iota(y)\rVert \le e^{c + \val_\ell(C_1)\log\ell}\).
\end{problem}

For the denominator clause of regularity, the proof uses the ramification and Selmer bounds of Theorem~\ref{thm:twisting} for the main summands of the lifts of good forms and of the rank-one lifts attached to main pairs, applies Theorem~\ref{dp-thm:denominators} to the image of the good-form lift and to the line of the rank-one lift, and adds \(\val_\ell(C_1)\log\ell\) to the logarithmic allowance for the rank-one pairing (as in Problem~\ref{pr:class-number}), where \(C_1\) is the number of \(U(1)\) double cosets at the chosen level.  For every finite set of primes the sum of the resulting logarithmic allowances is at most \(B\) (Problem~\ref{dp-pr:budget-sum}).

This chapter has deduced the small twists from the polynomial input and the lemma on roots of unity, and the small ratios of a regular family from the twisting theorem.  Taken as given: the polynomial input, the twisting theorem (Theorem~\ref{thm:twisting}), Theorem~\ref{dp-thm:denominators} and the existence of main summands used in the denominator clause, all read from the proof.  Chapter~\ref{tc-chap} uses Theorem~\ref{thm:regular}.

\paragraph{In the Lean formalization.}
In the Lean proof \cite{GL} (\nolinkurl{grh/src/}, namespace \texttt{New.Selmer} unless stated), the input of Section~\ref{sec:twisting} is a specialization of \texttt{exists\_twistPolynomials}, its objects are \texttt{SelmerDatum} and \texttt{StickelbergerSystem}, and the length bound is \texttt{SelmerDatum.TwistedLengthLe}.  Theorem~\ref{thm:small} is \texttt{exists\_anticyclotomic\_twist}, with conclusions \texttt{IsSmallTwist}; Theorem~\ref{thm:twisting} is \texttt{exists\_mFoldTwist\_hasDivisibleSelmerLengthLe}; Theorem~\ref{thm:regular} is \texttt{New.exists\_isRegularFamily}, and its clause on denominators is \texttt{exists\_mFoldTwist\_hasDenominatorsLe}.

\section*{Hints}

\subsection*{Section~\ref{sec:twisting}}
\begin{description}
\item[\ref{pr:psi-power}] For the last claim, the \(\lg_{w'}\) with \(w' \ne w\) are unramified at \(w\) because \(Q\) is auxiliary.
\item[\ref{pr:inertia-sum}] \(\lg_w\) is trivial on wild inertia, and a generator of tame inertia at \(w\) maps to a generator of \(\Z/\ell^{h_1}\); so \(\chi_\zeta(s_w) = \zeta_w^a\) with \(a\) prime to \(\ell\), and \(\zeta_w^{an}\), \(\zeta_w^n\) have the same order.  At \(\bar w\) use \(\lg_w(g^c) = -\lg_w(g)\).  Then Problem~\ref{w1-pr:powers-far}.
\item[\ref{pr:twist-from-roots}] Choose \(s_0\) as in the input, then a rational \(\rho > 0\) with \(s_0\rho \le \varepsilon/2\), then \(h_0\) by Problem~\ref{w1-pr:powers} with the bounds \(\bd + 1 \ge 0\) (note \(0 \le \val(F) \le \bd + s_0\rho\)) and \(\varepsilon/2\), then a sufficiently large \(h_1 \ge h_0\) and the auxiliary data at level \(\ell^{h_1}\).  Use Problem~\ref{pr:psi-power}; for the ramification, Problem~\ref{w1-pr:powers-far}; for the last sum, Problem~\ref{pr:inertia-sum} and \(2\,\#Q \le s_0\).
\item[\ref{pr:vl-log}] Factor \(n\); for the last claim take \(n = \mathcal N\).
\end{description}

\subsection*{Section~\ref{sec:regular}}
\begin{description}
\item[\ref{pr:height}] Which multiplier leaves the local inequalities unchanged?
\item[\ref{pr:coprime}] A place of \(Q\) has prime norm and does not lie above \(\ell\).
\item[\ref{pr:loss-sum}] Problem~\ref{pr:vl-log}, and count the factors of \(\mathrm{Int}_b\).
\item[\ref{pr:budget}] Write \(a = xb\).
\item[\ref{pr:assembly}] Problems~\ref{pr:coprime}, \ref{pr:inverse} and~\ref{pr:height}.
\item[\ref{pr:class-number}] \(\lVert\iota(C_1)\rVert = \ell^{-\val_\ell(C_1)}\).
\end{description}

\begin{thebibliography}{GL}
\bibitem[GL]{GL} L.~Sernau, \emph{GRH for Dirichlet \(L\)-functions, in Lean~4}, a shorter rewriting of the Lean proof, text by Claude, \url{https://github.com/lsernau/grh-lean}, commit \texttt{72fead1}, readable edition \texttt{grh/}.
\bibitem[Ne]{Ne} J.~Neukirch, \emph{Algebraic Number Theory}, Grundlehren Math.\ Wiss.\ 322, Springer, 1999.
\bibitem[Wa]{Wa} L.~C.~Washington, \emph{Introduction to Cyclotomic Fields}, 2nd ed., Grad.\ Texts in Math.\ 83, Springer, 1997.
\end{thebibliography}

\end{document}
